% Notes in the Laboratory house style: lecture, revision or derivation notes
% with definitions, worked examples and code. Replace the pendulum with the
% topic and keep the conventions. Build with build.py, which reads latexmkrc.
\documentclass[11pt]{article}
\usepackage{laboratory}

% The kind of document or the module code, set above the title; delete to omit.
\kicker{Lecture notes}
\title{Notes on the simple pendulum}
\author{Author Name}
\date{\today}
% The repository, notebook or lecture behind the notes; delete to omit.
\source{analysis/pendulum.py}

\begin{document}

\maketitle

\tableofcontents

\section{Equation of motion}

% Open each section with what it establishes, then define, derive and apply.
\begin{definition}
  A simple pendulum is a point mass $m$ on a massless, inextensible string of length $L$, swinging in a vertical plane under uniform gravity $g$ without friction.
  \label{def:pendulum}
\end{definition}

The bob moves on a circle of radius $L$, so its tangential acceleration is $L \ddot\theta$, and the only tangential force is the component of its weight, $-mg \sin\theta$. Newton's second law along the tangent gives
\begin{align}
  m L \frac{\dd^2\theta}{\dd t^2} &= -m g \sin\theta , \notag \\
  \frac{\dd^2\theta}{\dd t^2} &= -\frac{g}{L} \sin\theta .
  \label{eq:motion}
\end{align}
The mass cancels, so every pendulum of \cref{def:pendulum} with the same length swings alike~\cite{nelson1986}.

\section{Small oscillations}

For $|\theta| \ll 1$, $\sin\theta \approx \theta$ and \cref{eq:motion} becomes simple harmonic motion with angular frequency $\omega_0 = \sqrt{g/L}$ and period
\begin{equation}
  T_0 = 2\pi \sqrt{L/g} .
  \label{eq:period}
\end{equation}

\begin{example}
  A pendulum of length \qty{1.000}{\metre} where $g = \qty{9.81}{\metre\per\second\squared}$ has $T_0 = \qty{2.006}{\second}$, which is why a pendulum near a metre long beats seconds.
\end{example}

\Cref{fig:period} compares \cref{eq:period} with periods timed at nine lengths; the residuals in panel (b) stay within \qty{2}{\milli\second}, consistent with the timing error.

\begin{figure}[tbp]
  \centering
  \includegraphics[width=\linewidth]{figures/pendulum}
  \caption{\lead{The small-angle model fits every length.} Measured period against length~\marker{seriesA} with the small-angle model~\linekey{seriesC} (a), and residuals from the fit (b). Error bars are one standard error of a period; the band in (b) is one standard error of the fit.}
  \label{fig:period}
\end{figure}

\section{Finite amplitude}

Released from rest at $\theta_0$, the pendulum's period is an elliptic integral, which expands as~\cite{nelson1986}
\begin{equation}
  T = T_0 \left( 1 + \frac{\theta_0^2}{16} + \frac{11 \theta_0^4}{3072} + \cdots \right) .
  \label{eq:finite}
\end{equation}
The first correction is \qty{0.2}{\percent} at \qty{10}{\degree} and \qty{1.7}{\percent} at \qty{30}{\degree}, so a clock pendulum keeps time only if its amplitude is steady. \Cref{lst:fit} fits \cref{eq:period} to measured periods by least squares~\cite{taylor1997}.

\begin{lstlisting}[caption={Least-squares fit of $T = a \sqrt{L}$ and the resulting $g$.}, label=lst:fit]
def fit_g(lengths, periods):
    roots = np.sqrt(lengths)
    a = np.sum(periods * roots) / np.sum(lengths)
    return (2 * np.pi / a) ** 2
\end{lstlisting}

\bibliographystyle{apsrev4-2}
\bibliography{refs}

\end{document}
